4.- 
Utilizando doble contención \\
$\subseteq$ $\mid$

P.D que $L\subseteq L^{+}$

Por definción de cerradura inductiva dada en clase:
 
$L^{+}=\bigcup_{n\ge1}L^{n}=L^{1}\cup L^{2}\cup\dots$

Para $n=1$ aplicando la definicón de potencia:
$L^{1}=L \cdot L^{0}=L \cdot \{\lambda\}=L$

Dado que $L^{1}=L$, se cumple trivialmente que 
$$L\subseteq L^{+}$$

$\supseteq$ $\mid$

Como $L^{+}=\bigcup_{n\ge1}L^{n}$ para demostrar que $L^{+}\subseteq L$ hacemos inducción sobre $n\ge1$.

-Caso Base: $(n=1)$:
$$L^{1}=L\cdot L^{0} \overset{\text{def}}{=} L\{\lambda\}=L$$

Dado que $L^{1}=L$, se cumple trivialmente que
$$L \subseteq L$$

-Hipótesis de Inducción:
Supongamos que se cumple para k, es decir:
$$L^{k}\subseteq L$$

-Paso Inductivo:
P.D que $L^{k+1}\subseteq L$ \\
Por definicón de potencia:
$$L^{k+1}=L^{k}\cdot L$$

Por Hipotesis de Inducción sabemos que:
$$L^{k}\subseteq L$$
Entonces:
$$L^{k+1}=L^{k}\cdot L \subseteq L\cdot L$$
Por H.I del problema $(L^{2} \subseteq L)$
Asi:
$$L^{k+1} = L^{k}\cdot L$$
$$L^{k} \cdot L \subseteq L^{2} \subseteq L$$
Entonces:
$$L^{k+1} \subseteq L$$

Por lo tanto
$$L^{+}=\bigcup_{n\ge1}L^{n}\subseteq L$$

Al haber demostrado ambas contenciones, concluimos que:
$$L^{+}=L$$